\section{Einleitung}
\subsection{Problemstellung}
In vielen Unternehmen werden Bewerbungen noch über unstrukturierte Kanäle wie E-Mail oder manuelle Ablagen verwaltet. 
Dadurch fehlt häufig eine zentrale Übersicht über den Status einzelner Bewerbungen, was zu erhöhtem Verwaltungsaufwand und mangelnder Transparenz im Recruitingprozess führen kann. 
Gleichzeitig müssen bei der Verarbeitung personenbezogener Bewerberdaten rechtliche Anforderungen wie die Datenschutz-Grundverordnung sowie Gleichbehandlungsgrundsätze gemäß Allgemeines Gleichbehandlungsgesetz eingehalten werden. 

\subsection{Projektziel}
Ziel dieses Projektes ist die Konzeption und Entwicklung eines webbasierten Bewerbermanagementportals (Applicant Tracking System, ATS), das den Bewerbungsprozess digital, strukturiert und nachvollziehbar abbildet. 
Das System soll Bewerber:innen ermöglichen, offene Stellen einzusehen und Bewerbungen inklusive Dokumenten-Upload online einzureichen. 
Für Recruiter:innen wird ein geschützter Verwaltungsbereich bereitgestellt, in dem Stellenanzeigen erstellt, Bewerbungen eingesehen und bearbeitet werden können.
Neben der funktionalen Umsetzung werden rechtliche Anforderungen, insbesondere die Vorgaben der Datenschutz-Grundverordnung (DSGVO) sowie des Allgemeines Gleichbehandlungsgesetz (AGG), konzeptionell berücksichtigt.
Das Ergebnis des Projektes ist ein funktionsfähiger Prototyp (MVP), der die zentralen Kernprozesse eines digitalen Bewerbermanagementsystems abbildet.
